\chapter{Experiments and Results}\label{ch:experiments}
\thispagestyle{pagestyle}

\section{Experimental Setup}
\label{sec:setup}

All training experiments were conducted on a workstation running
Windows~11 with the Unity Editor invoked via WSL2 for the Python
process. The environment is summarised in \cref{table:setup}.

\begin{table}[H]
\caption{Experimental environment.}
\label{table:setup}
\begin{tabular}{|p{4cm}|p{7cm}|}
  \hline
  \textbf{Component} & \textbf{Specification} \\
  \hline
  CPU              & Intel Core i7-12700H (14 cores) \\
  \hline
  RAM              & 32 GB DDR5 \\
  \hline
  GPU              & NVIDIA RTX 3070 Ti (8 GB) \\
  \hline
  OS               & Windows 11 / WSL2 Ubuntu 22.04 \\
  \hline
  Unity            & 2022.3 LTS \\
  \hline
  Python           & 3.10.12 \\
  \hline
  Stable-Baselines3 & 2.3.0 \\
  \hline
  PyTorch          & 2.1.0 \\
  \hline
  Gymnasium        & 0.29.1 \\
  \hline
  MQTT broker      & Mosquitto 2.0.18 \\
  \hline
\end{tabular}
\end{table}

The Mosquitto broker ran on the host at \texttt{127.0.0.1:1883}. Unity
physics ran at the default fixed time-step of 0.02~s (50~Hz) with
\texttt{Time.timeScale = 1} (real time). The Python process executed in
a WSL2 terminal and communicated with Unity exclusively over the loopback
interface.

\section{Training Convergence}
\label{sec:convergence}

A baseline training run of 100{,}000 timesteps was conducted using the
default hyperparameters from \cref{table:ppo-params}. The episode reward
mean, computed over a rolling 10-episode window as reported by
Stable-Baselines3, is the primary convergence indicator.

At the start of training the agent explores nearly uniformly, producing
episode rewards in the range $[-5, +10]$. By approximately 20{,}000
timesteps the agent consistently avoids the inner wall and begins
accumulating progress rewards, with the rolling mean rising to
$[20, 40]$. The half-turn bonus of $+25$ is reliably triggered from
around 30{,}000 timesteps. By 80{,}000 timesteps the agent completes
full laps in the majority of episodes and the rolling mean stabilises
in the $[60, 80]$ range. \Cref{table:convergence} summarises milestones
across three independent seeds.

\begin{table}[H]
\caption{Training convergence milestones (mean over three seeds).}
\label{table:convergence}
\begin{tabular}{|p{6cm}|p{3cm}|p{2.5cm}|}
  \hline
  \textbf{Milestone} & \textbf{Timestep} & \textbf{Std. dev.} \\
  \hline
  First collision-free episode       & 12{,}400  & $\pm$2100 \\
  \hline
  Half-turn bonus reliably triggered & 31{,}000  & $\pm$3800 \\
  \hline
  First full lap completed           & 48{,}500  & $\pm$5200 \\
  \hline
  Rolling mean $\geq 60$             & 76{,}000  & $\pm$6900 \\
  \hline
\end{tabular}
\end{table}

\section{Reward Shaping Ablation}
\label{sec:ablation}

To evaluate the contribution of the half-turn milestone bonus, a second
training run was performed with $c_h = 0$ (all other parameters
unchanged). Without the half-turn bonus, the first complete lap was
delayed by approximately 18{,}000 timesteps on average and the rolling
mean reward at 100{,}000 steps was approximately 12 points lower. This
confirms that the intermediate milestone provides a useful progress
signal early in training.

A further ablation with reduced collision penalty ($c_c = -1$ instead
of $-5$) produced an agent that completed laps but drove closer to the
inner wall, indicating that the stronger penalty is necessary to
encourage a safety margin.

\section{Inference Evaluation}
\label{sec:inference-eval}

The \texttt{best\_model.zip} checkpoint was evaluated over 50 episodes
using \texttt{ai\_driver.py}. Results are summarised in
\cref{table:inference}.

\begin{table}[H]
\caption{Inference evaluation over 50 episodes.}
\label{table:inference}
\begin{tabular}{|p{7cm}|p{4cm}|}
  \hline
  \textbf{Metric} & \textbf{Value} \\
  \hline
  Lap completion rate         & 82\% (41/50) \\
  \hline
  Mean episode reward         & 63.4 \\
  \hline
  Mean lap duration           & 38.2 s \\
  \hline
  Mean speed (completed laps) & 47.6 km/h \\
  \hline
  Collisions per episode      & 0.6 \\
  \hline
  Timeout episodes            & 9 (18\%) \\
  \hline
\end{tabular}
\end{table}

The 18\% of episodes that ended in timeout typically corresponded to
the agent entering a slow-speed oscillation near an obstacle, suggesting
that the current reward shaping does not sufficiently penalise stalling.
This is identified as a direction for future improvement.

\section{Health Monitoring Validation}
\label{sec:health-validation}

The health monitoring system was validated by inducing controlled
failures:
\begin{itemize}[topsep=1.15pt,itemsep=1.15pt]
  \item \textbf{Observation timeout.} The Unity simulation was paused
    mid-episode. The health evaluator correctly transitioned to
    \textit{Critical} after 5 seconds with the message ``No observation
    received for 5 s''.
  \item \textbf{Python crash.} The Python process was killed with
    SIGKILL. The Editor detected the non-zero exit code and displayed
    ``Python process exited unexpectedly (exit code 137)''.
  \item \textbf{MQTT disconnect.} Mosquitto was stopped. The
    EditorMqttListener raised a connection-lost event within
    approximately 35 seconds (one keep-alive period) and the banner
    transitioned to \textit{Disconnected}.
\end{itemize}
All three failure modes were correctly detected and displayed without
manual intervention.

\section{End-to-End SDV Integration Test}
\label{sec:e2e-test}

The full deployment chain was validated on the QEMU-hosted Leda image.
With all three containers running, the following checks were performed:
\begin{enumerate}[leftmargin=2cm,topsep=1.15pt,itemsep=1.15pt,partopsep=1.15pt,parsep=1.15pt]
  \item A Kuksa \texttt{GetValue} call for \texttt{Vehicle.Speed}
    returned the current simulator speed within 100~ms of it being
    published to MQTT, confirming correct feeder operation.
  \item When the Python agent applied maximum throttle for 5 seconds,
    the Velocitas application published an overspeed alert to
    \texttt{leda/command/alert}, confirming end-to-end signal
    propagation from the simulation through the VSS layer to the safety
    monitor.
  \item The alert was received by a test \texttt{mosquitto\_sub}
    subscriber on the host machine, confirming that the outbound MQTT
    path from the Leda image to the host network was functional.
\end{enumerate}
